\section*{Zusammenfassung}
\addcontentsline{toc}{section}{Zusammenfassung}

In diesem Projekt wurde untersucht, wie gut vortrainierte neuronale Netze zur dreidimensionalen Objekterkennung auf eine fest installierte Zwei-Sensor-LiDAR-Anlage in einer Tiefgarage angepasst werden können. Die Zielklassen sind \textbf{Person}, \textbf{Fahrrad} und \textbf{Auto}. Verglichen wurden die beiden Einzelsensoransichten \texttt{os0} und \texttt{os1} sowie die aus beiden Sensoren erzeugte fusionierte Punktwolke \texttt{merged}.

Als Hauptmodell wurde PointPillars eingesetzt und ausgehend von einem auf KITTI vortrainierten Drei-Klassen-Checkpoint für 50 Epochen auf den Projektdaten finetuned. Wegen des starken Klassenungleichgewichts wurde eine ausschließlich aus Trainingsdaten erzeugte GT-Sampling-Datengrundlage verwendet. Die abschließende Bewertung folgt einer vor der Ergebnissichtung eingefrorenen gepaarten 3-Fold-Cross-Validation. In jedem Fold werden drei vollständige Experimente, je eines für Auto, Fahrrad und Person, als ungesehene Testaufnahmen zurückgehalten.

Die Anpassung an die Zielszene ist der entscheidende Schritt. In einem nicht kontrollierten Größenordnungsvergleich mit einer Vorgängerarbeit auf denselben Aufnahmen erreichen ausschließlich vortrainierte Modelle einen mAP von höchstens 0{,}058, das hier finetunete Modell 0{,}798. Der Abstand von 0{,}740 mAP ist je nach Ansicht sechs- bis sechzehnmal größer als der Abstand zwischen den beiden untersuchten Architekturen.

PointPillars erreicht mit der fusionierten Ansicht in jedem Fold den höchsten Gesamt-mAP. Der Fold-Mittelwert beträgt \(0{,}798 \pm 0{,}037\) für \texttt{merged}, \(0{,}777 \pm 0{,}018\) für \texttt{os0} und \(0{,}770 \pm 0{,}027\) für \texttt{os1}. Dieser Gesamtwert bewertet allerdings auch die Wiedererkennung der zahlreichen geparkten Fahrzeuge. Beschränkt auf die bewegten Objekte wächst der Vorsprung gegenüber der besseren Einzelansicht von 0{,}021 auf 0{,}040 in der offiziellen und 0{,}069 in der labelbasierten Auswertung. Der Fusionsvorteil ist moderat, aber über alle Folds konsistent und zeigt sich vor allem bei der präzisen Lokalisierung von Personen und Fahrrädern -- in beiden Auswertungsvarianten. Beim reinen Auffinden von Objekten ist die Fusion dagegen nicht überlegen: Auf den einzelnen Testexperimenten erreicht die Einzelansicht \texttt{os0} bei AP\textsubscript{30} sowohl den höchsten Mittelwert als auch die geringste Streuung. Zusätzlich wurde CenterPoint als Alternativarchitektur herangezogen; sie bestätigt die Rangfolge der Ansichten, bleibt aber in allen Ansichten hinter PointPillars zurück.

Die fusionierte Punktwolke benötigt auf der Evaluations-GPU ungefähr 85\,ms pro Modellinferenz, gegenüber 45\,ms für \texttt{os0} und 43\,ms für \texttt{os1}. Diese Messung belegt die 10-Hz-Fähigkeit der Modellinferenz, nicht jedoch die Echtzeitfähigkeit des vollständigen Fusionssystems. Die vorhandene Offline-Merge-Pipeline benötigt im Median 2{,}25\,s pro Frame; ein optimierter Online-Merge mit einmalig bestimmter Extrinsik wurde noch nicht gemessen.
